% ===================================================================
\section{Dynamic and Probabilistic Layer}
\label{ch:dynamic}
% ===================================================================

\subsection{Overview}

This chapter covers two connected aspects of FRFP dynamics:

\begin{enumerate}
  \item \textbf{Epistemic state and credence degradation} --- how tacit knowledge
        degrades over time when not refreshed by human judgment.
  \item \textbf{Probabilistic stopping times and survival analysis} --- the
        stochastic model of pipeline resolution, connecting to the $\runToOmega$
        bijection.
\end{enumerate}

% -------------------------------------------------------------------
\subsection{Epistemic State}
% -------------------------------------------------------------------

\begin{definition}[Tacit state]
A tacit state records the current quality of human-grounded knowledge in a
pipeline execution. It is represented as a record with fields including a credence
value and a degradation parameter.
\end{definition}

\begin{definition}[Credence]
The credence of a tacit state is a rational number:
\[
  c \in [0, 1] \subset \mathbb{Q}
\]
\end{definition}

\begin{correction}[Lean Correction P-07]
The original paper uses credence notation without stating bounds explicitly.
The Lean formalization enforces bounds by construction:

\begin{lstlisting}
structure Credence where
  value  : Rat
  nonneg : 0 ≤ value
  le_one : value ≤ 1
\end{lstlisting}

\leanref{Frfp.Core.InstitutionalLayer.Credence}

Without these structural bounds, any arithmetic on credence could produce values
outside $[0,1]$, making the tacit degradation theorem unprovable.
\end{correction}

% -------------------------------------------------------------------
\subsection{Degradation}
% -------------------------------------------------------------------

\begin{definition}[Degradation operator $\delta$]
The degradation operator models the decay of tacit correctness over time when
not refreshed by human evaluation:
\[
  \delta : \mathrm{TacitState} \to \mathrm{TacitState}
\]
\end{definition}

\begin{definition}[Degradation formula --- corrected]
Given credence $c \in [0,1]$ and degradation rate $d \geq 0$, after $t$ steps
the credence is:
\[
  c'(t) = \max\!\Big(0,\ \min\!\Big(1,\ c - d \cdot t\Big)\Big)
\]
\end{definition}

\begin{correction}[Lean Correction P-07 (cont.)]
The original paper writes $c'(t) = c - d \cdot t$ without clamping. The Lean
formalization requires explicit clamping to keep credence within $[0,1]$; without
it, the non-negativity proof obligation fails.
\end{correction}

\begin{theorem}[Degradation monotonicity]
Credence is monotonically non-increasing under repeated degradation:
\[
  t_1 \leq t_2 \Rightarrow c'(t_2) \leq c'(t_1)
\]
\leanref{Frfp.Core.DynamicLayer.degradation\_monotone}
\end{theorem}

% -------------------------------------------------------------------
\subsection{Entropy Drift}
% -------------------------------------------------------------------

\begin{definition}[Entropy drift]
The entropy of the tacit component of a pipeline increases monotonically when no
human evaluation step occurs between AI steps:
\[
  \Delta H = H(t_2) - H(t_1) \geq 0 \quad
  \text{for } t_1 \leq t_2 \text{ with no HFD/TE between them}
\]
\end{definition}

\begin{theorem}[Entropy drift monotonicity]
Under continuous AI-only execution (no $\TE$ or $\HFD$ steps), the entropy of
the tacit component is non-decreasing.
\leanref{Frfp.Core.DynamicLayer.entropy\_drift\_monotone}
\end{theorem}

Entropy drift formalizes the FRFP claim that purely AI-driven pipelines without
human checkpoints progressively lose tacit validity. The $\TE$ primitive resets
the drift by incorporating human evaluation.

% -------------------------------------------------------------------
\subsection{Stopping Times}
% -------------------------------------------------------------------

\begin{definition}[Stopping time]
A stopping time $\tau$ for pipeline $P$ is a measurable random variable:
\[
  \tau : \Omega \to \mathbb{N} \cup \{\infty\}
\]
such that for all $n \in \mathbb{N}$, the event $\{\tau \leq n\}$ is measurable
with respect to the filtration generated by the first $n$ steps of the run.

\leanref{Frfp.Core.Probability.StoppingTime}
\end{definition}

\begin{definition}[Hazard rate]
The hazard rate $h(n)$ is the conditional probability that the pipeline resolves
at step $n$, given it has not resolved before step $n$:
\[
  h(n) = \mathbb{P}(\tau = n \mid \tau \geq n)
       = \frac{S(n-1) - S(n)}{S(n-1)}
\]
\end{definition}

\begin{definition}[Survival function]
The survival function $S(n)$ is the probability that the pipeline has not yet
resolved by step $n$:
\[
  S(n) = \mathbb{P}(\tau > n)
\]
\end{definition}

% -------------------------------------------------------------------
\subsection{Five Proven Probability Theorems}
% -------------------------------------------------------------------

The following five theorems were axiomatized in an earlier version of the
formalization (March 2026) and subsequently proved using Lean's Mathlib library
(April 2026). They are now machine-checked theorems, not assumptions.

\begin{theorem}[Survival product formula]
\[
  S(N) = \prod_{k=0}^{N} (1 - h(k))
\]
\leanref{Frfp.Core.ProbabilityProven.survival\_product\_formula}
\end{theorem}
\begin{proof}
By induction on $N$, using the definition of conditional probability and the
product rule for independent events. \qed
\end{proof}

\begin{theorem}[Hazard--survival relation]
\[
  h(n) = \frac{S(n-1) - S(n)}{S(n-1)}
\]
\leanref{Frfp.Core.ProbabilityProven.hazard\_survival\_relation}
\end{theorem}

\begin{theorem}[Survival monotonicity]
\[
  N \leq M \Rightarrow S(M) \leq S(N)
\]
\leanref{Frfp.Core.ProbabilityProven.survival\_monotone}
\end{theorem}

\begin{theorem}[Geometric survival]
When the hazard rate is constant, $h(n) = h$ for all $n$:
\[
  S(n) = (1-h)^{n+1}
\]
\leanref{Frfp.Core.ProbabilityProven.geometric\_survival}
\end{theorem}

\begin{theorem}[Infinite product limit]
When $h(k) > 0$ for all $k$:
\[
  \prod_{k=0}^{\infty} (1 - h(k)) = 0
\]
\leanref{Frfp.Core.ProbabilityProven.infinite\_product\_limit}
\end{theorem}

\begin{note}
The transition of these five theorems from axioms to machine-checked results
demonstrates that the FRFP probabilistic foundations are sound.
\end{note}

% -------------------------------------------------------------------
\subsection{Safe Horizon}
% -------------------------------------------------------------------

\begin{definition}[Safe horizon]
Given a degradation rate $d$ and an acceptability threshold $c_{\min}$, the safe
horizon $N^*(c, d)$ is the latest step at which the tacit state still has
acceptable credence:
\[
  N^*(c, d) = \max\{n \in \mathbb{N} \mid c'(n) \geq c_{\min}\}
            = \left\lfloor \frac{c - c_{\min}}{d} \right\rfloor
\]
\end{definition}

\begin{theorem}[Safe horizon characterization]
The safe horizon satisfies:
\[
  c'(N^*) \geq c_{\min} \wedge c'(N^*+1) < c_{\min}
\]
\leanref{Frfp.Core.DynamicLayer.safe\_horizon\_characterization}
\end{theorem}

The safe horizon determines when the pipeline \emph{must} invoke a $\TE$ or
$\HFD$ step --- i.e., when human evaluation is no longer optional.

% -------------------------------------------------------------------
\subsection{Tacit Irreversibility}
% -------------------------------------------------------------------

\begin{theorem}[Tacit irreversibility]
Once a pipeline transitions from explicit to tacit space (via $\RB$), this
transition cannot be reversed:
\[
  \nexists\ \text{sequence of primitives from } \Tobj \text{ to } \Eobj
\]
\leanref{Frfp.Core.TacitDependence}
\end{theorem}

This is a restatement of Theorem~1.7 (No-Return Principle) in dynamic terms:
once human evaluation is required, machine processing alone cannot restore
tacit validity.
